\documentclass[10pt,a4paper]{amsart}
\usepackage[margin=19mm]{geometry}
\usepackage{amsmath,amssymb,mathtools}
\usepackage[scr=rsfs,cal=euler]{mathalfa}
\usepackage{xfrac}
\usepackage[unicode,hidelinks,bookmarks=false]{hyperref}
\newcommand{\ie}{i.e.\ }
\newcommand{\cf}{cf.\ }
\newcommand{\resp}{respectively\ }
\newcommand{\Subsection}{\subsection}
\providecommand{\nobreakdash}{\nobreak\hbox{-}}
\setlength{\emergencystretch}{2em}
\multlinegap0pt
\usepackage[shortlabels]{enumitem}
\usepackage{amssymb}
\usepackage{mathtools}
\usepackage{float}
\usepackage{algorithm}\usepackage{algpseudocode}
\newtheorem{theorem}{Theorem}
\title{Repetition in Permutation Wordle}
\author{Aurora Hiveley}
\date{Source: arXiv:2601.05971v1 (CC BY 4.0)}
\begin{document}
\maketitle

\noindent\emph{Source and licence.} Repetition in Permutation Wordle. Version arXiv:2601.05971v1 (CC BY 4.0), distributed by arXiv under the Creative Commons Attribution 4.0 International licence, http://creativecommons.org/licenses/by/4.0/, as recorded in the arXiv licence metadata for this exact version and shown on the abstract page https://arxiv.org/abs/2601.05971v1. Full licence notice: \texttt{licenses/arxiv-2601.05971-CC-BY-4.0.txt}.\par\smallskip

\noindent\textit{Editorial scope.} Counted unit: the article's theorem theorem (source0.tex lines 244--246). The article's complete original proof of it is delivered with it (source0.tex lines 248--267, one \texttt{proof} environment of 20 lines). No mathematics is written, repaired or re-derived: the statement and the proof are the article's own lines, in the article's own notation, and no retained line is substituted or re-flowed.  No further source-local context is needed: every cross-reference of every retained line resolves inside the two delivered spans.  Nothing was omitted from any retained span, so the package carries no omission record.  No retained line contains a citation, so the wrapper carries no bibliography and no bibliography entry.  No retained line contains a figure environment, an image-inclusion command or drawing code, so no image is delivered and the package declares no asset dependency.  Editorial formatting only, in addition: an AMS \texttt{amsart} preamble that restates the article's own declarations and macros used by the retained lines, namely the environments \texttt{theorem} on separate counters, together with the standard packages those lines need.  The retained environments print in this document as Theorem 1 in order of appearance, whereas the article prints the same items with the numbering of its own full document.  The complete native source of the article as distributed in the arXiv e-print for this exact version is bundled unchanged as \texttt{sources/192380/source0.tex}.
\begin{theorem}
    Let $\delta$ be a derangement of length $n$ with cycle type $[t_1, t_2, \dots, t_k]$. Let $\mu_i := |\{ j \mid t_j = t_i \}|$, or in other words $\mu_i$ is the multiplicity of the cycle length $t_i$ in the cycle type of $\delta$. If $\mu_i \geq 2$ for all $1 \leq i \leq n$, then any strategy with $S[n] = \delta$ will enter an infinite loop for at least one possible secret permutation $\omega$ of length $n$.
\end{theorem}

\begin{proof}
Let $\delta$ be a derangement with with cycle type $[t_1, t_2, \dots, t_k]$ such that $\mu_i \geq 2$ for all $1 \leq i \leq k$. Let $\delta = c_1 c_2 \dots c_k$ be the partition of $\delta$ into cycles where the length of $c_i$ is $t_i$ for each $i$. 
%Without loss of generality, assume that partitioning $\delta$ into $c_i$'s takes the form $\delta = (1, \dots, t_1)(t_1+1, \dots ,t_1 + t_2) \dots (n-t_k, \dots, n)$ so that $c_i = \left( 1 + \sum_{j=1}^{i-1} t_j, 2 + \sum_{j=1}^{i-1} t_j, \dots, \sum_{j=1}^i t_j \right)$. 
For a unique cycle length $t_i$, let $\mathcal{C}_i = \{c_i^1, c_i^2, \dots, c_i^{\mu_i} \}$ be the set of cycles in $\delta$ that each have length $t_i$. Then to construct an offending permutation $\omega$, map the elements of $c_i^j$ to the positions occupied by $c_i^{j+1}$ in $\delta$, in other words $\omega[c_i^j[a]] = c_i^{j+1}[a]$ for any $1 \leq a \leq t_i$. Repeat this mapping for all $1 \leq j \leq \mu_i$ (modulo $\mu_i$, of course, so that $c_i^{\mu_i}$'s elements map to the positions occupied by $c_i^1$) and for all unique cycle lengths $t_i$. Consider the example below where $\delta$ has length 10 and decomposes into two 2-cycles and two 3-cycles:
\[ 
\delta = [2,3,1,5,4,7,8,6,10,9] 
= (1,2,3)(4,5)(6,7)(8,9,10) 
\implies \omega = [8,9,10,6,7,4,5,1,2,3]
\]

Observe that a permutation wordle player using the inductive strategy such that $S[n] = \delta$ will enter an infinite loop while attempting to guess $\omega$. The player's first guess is always the trivial permutation $[1,2,\dots, n]$, and since $\delta$ was a derangement, each cycle $c_i$ in $\delta$ has length at least 2. Since $\mu_i \geq 2$ for all $i$, no entries of $\omega$ are shared with the trivial permutation. Then $\delta = S[n]$ generates the next guess as follows:
\[
\begin{array}{cc}   
\gamma_1 = \underbrace{1,2,\dots,t_1}_{c_1},\underbrace{t_1 + 1, \dots, t_1 + t_2}_{c_2}, \dots, \underbrace{n - t_k, \dots, n}_{c_k}  \\  
\gamma_2 = \underbrace{2,\dots,t_1,1}_{},\underbrace{t_1 + 2, \dots, t_1 + t_2, t_1+1}_{}, \dots, \underbrace{n - t_k + 1, \dots, n, n-t_k}_{}  \\
\end{array}
\]

Once again, none of the entries will be correct in $\gamma_2$ since each entry of $c_1$ was mapped to a location within the first $t_1$ locations, and the same for $c_2$, and so on. Since $\mathcal{J}_2 = \emptyset$, we generate $\gamma_3$ according to $\delta$ once again, but by the same logic as before, $\mathcal{J}_3 = \emptyset$. This process will repeat infinitely without ever locking in any correct entries, so we cannot guess the permutation $\omega$ in finitely many guesses using $S[n] = \delta$. 
\end{proof}

\end{document}
